\documentclass[a4paper,10pt]{report}\input{../0.Préambule/Préambule_cours_prof}\begin{document}

\pagestyle{empty}
\newpage
\section*{Statistiques - Calcul littéral \hspace{2cm} Sujet A}
\addcontentsline{toc}{section}{Statistiques - Calcul littéral Sujet A}

\noindent \textit{L'usage de la calculatrice est \textbf{interdite}.} \hfill \textit{Durée : 55 minutes}

\medskip
\noindent \textit{Ceci est un questionnaire a choix multiples (QCM). Pour chaque question, il y a une ou plusieurs bonnes réponses.}

\medskip
\noindent \textit{Une bonne réponse rapporte un point, un mauvaise réponse enlève un tiers de point et une absence de réponse ne rapporte ni n'enlève aucun point.}

\subsection*{Calcul littéral}
\addcontentsline{toc}{subsection}{Calcul littéral}

\bigskip
\textbf{Affirmation 1 : } : L'expression développée et simplifiée de $3(x+6)$ est :

\medskip
\begin{tabular}{p{3.8cm}p{3.8cm}p{3.8cm}p{3.8cm}}
\textbf{A. } $3x+6$&
\textbf{B. } $x+18$&
\textbf{C. } $3x+18$&
\textbf{D. } $18x$\\
\end{tabular}

\bigskip
\textbf{Affirmation 2 : } $12x-42$ est la forme développée et simplifiée de l'expression :

\medskip
\begin{tabular}{p{3.8cm}p{3.8cm}p{3.8cm}p{3.8cm}}
\textbf{A. } $-6(2x+7)$&
\textbf{B. } $6(2x-7)$&
\textbf{C. } $6(-2x-7)$&
\textbf{D. } $6(-2x+7)$\\
\end{tabular}

\bigskip
\textbf{Affirmation 3 : }  L'expression développée et simplifiée de $-4(2x-6)$ est :

\medskip
\begin{tabular}{p{3.8cm}p{3.8cm}p{3.8cm}p{3.8cm}}
\textbf{A. } $-8x+24$&
\textbf{B. } $-8x-24$&
\textbf{C. } $8x-24$&
\textbf{D. } $8x+24$\\
\end{tabular}

\bigskip
\textbf{Affirmation 4 : } L'expression développée de $4+3(2y-2)$ est :

\medskip
\begin{tabular}{p{3.8cm}p{3.8cm}p{3.8cm}p{3.8cm}}
\textbf{A. } $7+2y-2$ &
\textbf{B. } $7y$ &
\textbf{C. } $6y-2$ &
\textbf{D. } $14y-14$ \\
\end{tabular}

\bigskip
\textbf{Affirmation 5 : } Soit $5(8-2x) = ax +40$. La valeur de $a$ est :

\medskip
\begin{tabular}{p{3.8cm}p{3.8cm}p{3.8cm}p{3.8cm}}
\textbf{A. } $10$&
\textbf{B. } $0$&
\textbf{C. } $-10$&
\textbf{D. } $50$\\
\end{tabular}

\bigskip
\noindent Pour les affirmations \textbf{6}, \textbf{7} et \textbf{8}, on s'intéresse au programme suivant :

\begin{center}
\arrayrulecolor{black}
\renewcommand{\arraystretch}{1.2}
\begin{tabular}{|l|}
\hline
\edupythoninline{Choisis un nombre}\\
\edupythoninline{Soustrais 6 a ce nombre}\\
\edupythoninline{Multiplie le resultat par -2}\\
\edupythoninline{Ajoute le triple du nombre de depart}\\
\hline
\end{tabular}
\end{center}

\bigskip
\textbf{Affirmation 6 : } Si on choisit $5$ au début de ce programme, combien trouve-t-on en sortie ?

\medskip
\begin{tabular}{p{3.8cm}p{3.8cm}p{3.8cm}p{3.8cm}}
\textbf{A. } $17$&
\textbf{B. } $10$&
\textbf{C. } $8$&
\textbf{D. } $4$\\
\end{tabular}

\bigskip
\textbf{Affirmation 7 : } Si on choisit $-2$ au début de ce programme, combien trouve-t-on en sortie ?

\medskip
\begin{tabular}{p{3.8cm}p{3.8cm}p{3.8cm}p{3.8cm}}
\textbf{A. } $17$&
\textbf{B. } $10$&
\textbf{C. } $8$&
\textbf{D. } $4$\\
\end{tabular}

\bigskip
\textbf{Affirmation 8 : } Si on choisit $x$ au début de ce programme, quelle expression trouve-t-on en sortie ?

\medskip
\begin{tabular}{p{3.8cm}p{3.8cm}p{3.8cm}p{3.8cm}}
\textbf{A. } $x$&
\textbf{B. } $4$&
\textbf{C. } $2(x-6)+3x$&
\textbf{D. } $x+12$\\
\end{tabular}

\bigskip
\noindent Pour les affirmations \textbf{9} et \textbf{10}, on considère un rectangle de longueur $(x+6)$ et de largeur $x$

\bigskip
\textbf{Affirmation 9 : } Le périmètre de ce rectangle peut être donnée par l'expression :

\medskip
\begin{tabular}{p{3.8cm}p{3.8cm}p{3.8cm}p{3.8cm}}
\textbf{A. } $2 \times (2x+6)$&
\textbf{B. } $2x + 6 +2x$&
\textbf{C. } $x + 6 + x$&
\textbf{D. } $(x+6) \times x$\\
\end{tabular}

\bigskip
\textbf{Affirmation 10 : } L'aire de ce rectangle peut être donnée par l'expression :

\medskip
\begin{tabular}{p{3.8cm}p{3.8cm}p{3.8cm}p{3.8cm}}
\textbf{A. } $6+2x$&
\textbf{B. } $6x^2$&
\textbf{C. } $12x$&
\textbf{D. } $x^2+6x$\\
\end{tabular}

\newpage
\subsection*{Statistiques}
\addcontentsline{toc}{subsection}{Statistiques}

\noindent Pour les affirmations \textbf{11} à \textbf{14}, on utilisera l'énoncé ci-dessous.

\noindent Pour calculer la moyenne du devoir de mathématiques d'une classe de quatrième , on a la formule $$\dfrac{1+2+7+8+9+7+8+12+14+15+17+18+3+20+8+5+9+7+12+14+18+20}{22}$$

\bigskip
\textbf{Affirmation 11 : } Le nombre d'élève dans cette classe est :

\medskip
\begin{tabular}{p{3.8cm}p{3.8cm}p{3.8cm}p{3.8cm}}
\textbf{A. } $12$&
\textbf{B. } $22$&
\textbf{C. } $4,36$&
\textbf{D. } $35$\\
\end{tabular}

\bigskip
\textbf{Affirmation 12 : } la médiane des notes à ce devoir est :

\medskip
\begin{tabular}{p{3.8cm}p{3.8cm}p{3.8cm}p{3.8cm}}
\textbf{A. } $9$&
\textbf{B. } $10$&
\textbf{C. } $8$&
\textbf{D. } $17,5$\\
\end{tabular}

\bigskip
\textbf{Affirmation 13 : } On donne le tableau valeurs effectifs des notes obtenus lors de ce contrôle :

\begin{center}
\renewcommand{\arraystretch}{1.4}
\begin{tabularx}{15cm}{|>{\columncolor{gray!50}}p{3cm}|*{13}{>{\centering \arraybackslash}X|}}
\hline
Note & $0$ & $1$ & $2$ & $3$ & $7$ & $8$ & $9$ & $12$ & $14$ & $15$ & $17$ & $18$ & $20$\\\hline
Effectif &  &  &  &  &  & \bctrefle &  &  &  &  &  &  & \\\hline
\end{tabularx}
\end{center}

Quelle est la valeur dans la case \quad \bctrefle \quad ?

\medskip
\begin{tabular}{p{3.8cm}p{3.8cm}p{3.8cm}p{3.8cm}}
\textbf{A. } $1$&
\textbf{B. } $2$&
\textbf{C. } $3$&
\textbf{D. } $4$\\
\end{tabular}

\bigskip
\textbf{Affirmation 14 : } L'étendue de cette série de note est :

\medskip
\begin{tabular}{p{3.8cm}p{3.8cm}p{3.8cm}p{3.8cm}}
\textbf{A. } $10$&
\textbf{B. } $20$&
\textbf{C. } $19$&
\textbf{D. } $22$\\
\end{tabular}

\bigskip
\noindent Pour les affirmations \textbf{15} à \textbf{17}, on utilisera l'énoncé ci-dessous.

\noindent Dans un groupe de personnes, on considère le nombre de frères et sœurs. On relève les données statistiques dans le tableau suivant :

\begin{center}
\renewcommand{\arraystretch}{1.2}
\begin{tabularx}{\linewidth}{|>{\columncolor{gray!50}}p{4cm}|*{8}{>{\centering \arraybackslash}X|}}
\hline
Nombre de frères et sœurs & $0$ & $1$ & $2$ & $3$ & $4$ & $5$ & $6$ & $7$\\\hline
Effectif & $3$ & $6$ & $7$ & $9$ & $5$ & $2$ & $1$ & $1$\\\hline
\end{tabularx}
\end{center}

\bigskip
\textbf{Affirmation 15 : } L'effectif total de cette série est :

\medskip
\begin{tabular}{p{3.8cm}p{3.8cm}p{3.8cm}p{3.8cm}}
\textbf{A. } $7$&
\textbf{B. } $9$&
\textbf{C. } $34$&
\textbf{D. } $12$\\
\end{tabular}

\bigskip
\textbf{Affirmation 16 : } Combien de personnes ont au moins trois frères et sœurs ?

\medskip
\begin{tabular}{p{3.8cm}p{3.8cm}p{3.8cm}p{3.8cm}}
\textbf{A. } $18$&
\textbf{B. } $5$&
\textbf{C. } $3$&
\textbf{D. } $16$\\
\end{tabular}

\bigskip
\textbf{Affirmation 17 : } Le nombre moyen de frères et sœurs est :

\medskip
\begin{tabular}{p{3.8cm}p{3.8cm}p{3.8cm}p{3.8cm}}
\textbf{A. } $3,5$&
\textbf{B. } inférieur à $2$&
\textbf{C. } supérieur à $4$&
\textbf{D. } entre $2$ et $3$\\
\end{tabular}

\bigskip
\noindent Pour les question \textbf{18} à \textbf{20}, on s'intéresse à la série statistique suivante :
\begin{center}
$2$ – $14$ – $7$ – $11$ – $9$ – $9$ – $10$ – $10$ – $11$ – $4$
\end{center}

\bigskip
\textbf{Affirmation 18 : } La médiane vaut :

\medskip
\begin{tabular}{p{3.8cm}p{3.8cm}p{3.8cm}p{3.8cm}}
\textbf{A. } $9,5$&
\textbf{B. } $9$&
\textbf{C. } $2$&
\textbf{D. } $12$\\
\end{tabular}

\bigskip
\textbf{Affirmation 19 : } L'étendue vaut :

\medskip
\begin{tabular}{p{3.8cm}p{3.8cm}p{3.8cm}p{3.8cm}}
\textbf{A. } $9,5$&
\textbf{B. } $9$&
\textbf{C. } $2$&
\textbf{D. } $12$\\
\end{tabular}

\bigskip
\textbf{Affirmation 20 : } La moyenne est :

\medskip
\begin{tabular}{p{3.8cm}p{3.8cm}p{3.8cm}p{3.8cm}}
\textbf{A. } négative &
\textbf{B. } supérieure à $14$&
\textbf{C. } inférieur à $4$&
\textbf{D. } un nombre décimal\\
\end{tabular}

\newpage
\section*{Statistiques - Calcul littéral \hspace{2cm} Sujet B}
\addcontentsline{toc}{section}{Statistiques - Calcul littéral Sujet A}

\noindent \textit{L'usage de la calculatrice est \textbf{interdite}.} \hfill \textit{Durée : 55 minutes}

\medskip
\noindent \textit{Ceci est un questionnaire a choix multiples (QCM). Pour chaque question, il y a une ou plusieurs bonnes réponses.}

\medskip
\noindent \textit{Une bonne réponse rapporte un point, un mauvaise réponse enlève un tiers de point et une absence de réponse ne rapporte ni n'enlève aucun point.}

\subsection*{Calcul littéral}
\addcontentsline{toc}{subsection}{Calcul littéral}

\bigskip
\textbf{Affirmation 1 : } : L'expression développée et simplifiée de $3(x+6)$ est :

\medskip
\begin{tabular}{p{3.8cm}p{3.8cm}p{3.8cm}p{3.8cm}}
\textbf{A. } $3x+6$&
\textbf{B. } $3x+18$&
\textbf{C. } $x+18$&
\textbf{D. } $18x$\\
\end{tabular}

\bigskip
\textbf{Affirmation 2 : }  L'expression développée et simplifiée de $-4(2x-6)$ est :

\medskip
\begin{tabular}{p{3.8cm}p{3.8cm}p{3.8cm}p{3.8cm}}
\textbf{A. } $-8x+24$&
\textbf{B. } $8x-24$&
\textbf{C. } $-8x-24$&
\textbf{D. } $8x+24$\\
\end{tabular}

\bigskip
\textbf{Affirmation 3 : } $12x-42$ est la forme développée et simplifiée de l'expression :

\medskip
\begin{tabular}{p{3.8cm}p{3.8cm}p{3.8cm}p{3.8cm}}
\textbf{A. } $-6(2x+7)$&
\textbf{B. } $6(-2x-7)$&
\textbf{C. } $6(2x-7)$&
\textbf{D. } $6(-2x+7)$\\
\end{tabular}

\bigskip
\textbf{Affirmation 4 : } L'expression développée de $4+3(2y-2)$ est :

\medskip
\begin{tabular}{p{3.8cm}p{3.8cm}p{3.8cm}p{3.8cm}}
\textbf{A. } $7+2y-2$ &
\textbf{B. } $6y-2$ &
\textbf{C. } $7y$ &
\textbf{D. } $14y-14$ \\
\end{tabular}

\bigskip
\textbf{Affirmation 5 : } Soit $5(8-2x) = ax +40$. La valeur de $a$ est :

\medskip
\begin{tabular}{p{3.8cm}p{3.8cm}p{3.8cm}p{3.8cm}}
\textbf{A. } $10$&
\textbf{B. } $-10$&
\textbf{C. } $0$&
\textbf{D. } $50$\\
\end{tabular}

\bigskip
\noindent Pour les affirmations \textbf{6}, \textbf{7} et \textbf{8}, on s'intéresse au programme suivant :

\begin{center}
\arrayrulecolor{black}
\renewcommand{\arraystretch}{1.2}
\begin{tabular}{|l|}
\hline
\edupythoninline{Choisis un nombre}\\
\edupythoninline{Soustrais 6 a ce nombre}\\
\edupythoninline{Multiplie le resultat par -2}\\
\edupythoninline{Ajoute le triple du nombre de depart}\\
\hline
\end{tabular}
\end{center}

\bigskip
\textbf{Affirmation 6 : } Si on choisit $-2$ au début de ce programme, combien trouve-t-on en sortie ?

\medskip
\begin{tabular}{p{3.8cm}p{3.8cm}p{3.8cm}p{3.8cm}}
\textbf{A. } $17$&
\textbf{B. } $8$&
\textbf{C. } $10$&
\textbf{D. } $4$\\
\end{tabular}

\bigskip
\textbf{Affirmation 7 : } Si on choisit $5$ au début de ce programme, combien trouve-t-on en sortie ?

\medskip
\begin{tabular}{p{3.8cm}p{3.8cm}p{3.8cm}p{3.8cm}}
\textbf{A. } $17$&
\textbf{B. } $8$&
\textbf{C. } $10$&
\textbf{D. } $4$\\
\end{tabular}

\bigskip
\textbf{Affirmation 8 : } Si on choisit $x$ au début de ce programme, quelle expression trouve-t-on en sortie ?

\medskip
\begin{tabular}{p{3.8cm}p{3.8cm}p{3.8cm}p{3.8cm}}
\textbf{A. } $x$&
\textbf{B. } $2(x-6)+3x$&
\textbf{C. } $4$&
\textbf{D. } $x+12$\\
\end{tabular}

\bigskip
\noindent Pour les affirmations \textbf{9} et \textbf{10}, on considère un rectangle de longueur $(x+6)$ et de largeur $x$


\bigskip
\textbf{Affirmation 9 : } L'aire de ce rectangle peut être donnée par l'expression :

\medskip
\begin{tabular}{p{3.8cm}p{3.8cm}p{3.8cm}p{3.8cm}}
\textbf{A. } $6+2x$&
\textbf{B. } $12x$&
\textbf{C. } $6x^2$&
\textbf{D. } $x^2+6x$\\
\end{tabular}

\bigskip
\textbf{Affirmation 10 : } Le périmètre de ce rectangle peut être donnée par l'expression :

\medskip
\begin{tabular}{p{3.8cm}p{3.8cm}p{3.8cm}p{3.8cm}}
\textbf{A. } $2 \times (2x+6)$&
\textbf{B. } $x + 6 + x$&
\textbf{C. } $2x + 6 +2x$&
\textbf{D. } $(x+6) \times x$\\
\end{tabular}



\newpage
\subsection*{Statistiques}
\addcontentsline{toc}{subsection}{Statistiques}

\noindent Pour les affirmations \textbf{11} à \textbf{14}, on utilisera l'énoncé ci-dessous.

\noindent Pour calculer la moyenne du devoir de mathématiques d'une classe de quatrième , on a la formule $$\dfrac{1+2+7+8+9+7+8+12+14+15+17+18+3+20+8+5+9+7+12+14+18+20}{22}$$

\bigskip
\textbf{Affirmation 11 : } la médiane des notes à ce devoir est :

\medskip
\begin{tabular}{p{3.8cm}p{3.8cm}p{3.8cm}p{3.8cm}}
\textbf{A. } $9$&
\textbf{B. } $8$&
\textbf{C. } $10$&
\textbf{D. } $17,5$\\
\end{tabular}

\bigskip
\textbf{Affirmation 12 : } Le nombre d'élève dans cette classe est :

\medskip
\begin{tabular}{p{3.8cm}p{3.8cm}p{3.8cm}p{3.8cm}}
\textbf{A. } $12$&
\textbf{B. } $4,36$&
\textbf{C. } $22$&
\textbf{D. } $35$\\
\end{tabular}

\bigskip
\textbf{Affirmation 13 : } On donne le tableau valeurs effectifs des notes obtenus lors de ce contrôle :

\begin{center}
\renewcommand{\arraystretch}{1.4}
\begin{tabularx}{15cm}{|>{\columncolor{gray!50}}p{3cm}|*{13}{>{\centering \arraybackslash}X|}}
\hline
Note & $0$ & $1$ & $2$ & $3$ & $7$ & $8$ & $9$ & $12$ & $14$ & $15$ & $17$ & $18$ & $20$\\\hline
Effectif &  &  &  &  &  &  &  &  &  & \bctrefle &  &  & \\\hline
\end{tabularx}
\end{center}

Quelle est la valeur dans la case \quad \bctrefle \quad ?

\medskip
\begin{tabular}{p{3.8cm}p{3.8cm}p{3.8cm}p{3.8cm}}
\textbf{A. } $1$&
\textbf{B. } $3$&
\textbf{C. } $2$&
\textbf{D. } $4$\\
\end{tabular}

\bigskip
\textbf{Affirmation 14 : } L'étendue de cette série de note est :

\medskip
\begin{tabular}{p{3.8cm}p{3.8cm}p{3.8cm}p{3.8cm}}
\textbf{A. } $10$&
\textbf{B. } $19$&
\textbf{C. } $20$&
\textbf{D. } $22$\\
\end{tabular}

\bigskip
\noindent Pour les affirmations \textbf{15} à \textbf{17}, on utilisera l'énoncé ci-dessous.

\noindent Dans un groupe de personnes, on considère le nombre de frères et sœurs. On relève les données statistiques dans le tableau suivant :

\begin{center}
\renewcommand{\arraystretch}{1.2}
\begin{tabularx}{\linewidth}{|>{\columncolor{gray!50}}p{4cm}|*{8}{>{\centering \arraybackslash}X|}}
\hline
Nombre de frères et sœurs & $0$ & $1$ & $2$ & $3$ & $4$ & $5$ & $6$ & $7$\\\hline
Effectif & $3$ & $6$ & $7$ & $9$ & $5$ & $2$ & $1$ & $1$\\\hline
\end{tabularx}
\end{center}

\bigskip
\textbf{Affirmation 15 : } Combien de personnes ont au moins trois frères et sœurs ?

\medskip
\begin{tabular}{p{3.8cm}p{3.8cm}p{3.8cm}p{3.8cm}}
\textbf{A. } $18$&
\textbf{B. } $3$&
\textbf{C. } $5$&
\textbf{D. } $16$\\
\end{tabular}

\bigskip
\textbf{Affirmation 16 : } L'effectif total de cette série est :

\medskip
\begin{tabular}{p{3.8cm}p{3.8cm}p{3.8cm}p{3.8cm}}
\textbf{A. } $7$&
\textbf{B. } $34$&
\textbf{C. } $9$&
\textbf{D. } $12$\\
\end{tabular}

\bigskip
\textbf{Affirmation 17 : } Le nombre moyen de frères et sœurs est :

\medskip
\begin{tabular}{p{3.8cm}p{3.8cm}p{3.8cm}p{3.8cm}}
\textbf{A. } $3,5$&
\textbf{B. } supérieur à $4$&
\textbf{C. } inférieur à $2$&
\textbf{D. } entre $2$ et $3$\\
\end{tabular}

\bigskip
\noindent Pour les question \textbf{18} à \textbf{20}, on s'intéresse à la série statistique suivante :
\begin{center}
$2$ – $14$ – $7$ – $11$ – $9$ – $9$ – $10$ – $10$ – $11$ – $4$
\end{center}

\bigskip
\textbf{Affirmation 18 : } L'étendue vaut :

\medskip
\begin{tabular}{p{3.8cm}p{3.8cm}p{3.8cm}p{3.8cm}}
\textbf{A. } $9,5$&
\textbf{B. } $2$&
\textbf{C. } $9$&
\textbf{D. } $12$\\
\end{tabular}

\bigskip
\textbf{Affirmation 19 : } La moyenne est :

\medskip
\begin{tabular}{p{3.8cm}p{3.8cm}p{3.8cm}p{3.8cm}}
\textbf{A. } négative &
\textbf{B. } inférieur à $4$&
\textbf{C. } supérieure à $14$&
\textbf{D. } un nombre décimal\\
\end{tabular}

\bigskip
\textbf{Affirmation 20 : } La médiane vaut :

\medskip
\begin{tabular}{p{3.8cm}p{3.8cm}p{3.8cm}p{3.8cm}}
\textbf{A. } $9,5$&
\textbf{B. } $2$&
\textbf{C. } $9$&
\textbf{D. } $12$\\
\end{tabular}





\end{document}
